\cotefile{157-4}
\scanpage{333}
% @ top centre
\mspage{57}

% @ top, lines 1–2
\unsure{Or|Car} l'axiome $\mathrm{Loc}(8\mathrm{bis})$ implique que $f\in\mathrm{Fibf}_{W_0}$ se factorise en

% @ upper centre, triangle (α)
\diagram{\begin{tikzcd}[ampersand replacement=\&,column sep=small]
\& X \arrow[rr,"{i_{\infty}\deleted{\unsure{X|f}}}"] \arrow[dr,"f"'] \& \& \overline{X}=\unsure{\Phi|\mathscr L|\mathcal L}_{\infty}(X) \arrow[dl,"{p_{\infty}\deleted{\unsure{X|f}}}"]\\
\text{\unsure{($\alpha$)|($*$)}} \& \& S \& {,}
\end{tikzcd}}

% @ right of triangle
\margin{\deleted{\unsure{cat.}}\unsure{Cat.} de chemins infinis d'un côté, \unsure{relativisée} sur $S$ i.e. des couples $(x,c)$, avec $x\in\mathrm{Ob}\,X$ et $c:x\longrightarrow s$ un chemin infini dans $S$, de source $x$.}

% @ below triangle, lines 1–3
avec $p_\infty\in\mathrm{Parf}_{W_0}$ (comme on veut), et $\unsure{j_\infty|i_\infty}\in W_{0S}^{\unsure{u|n}}\deleted{\subset\unsure{W.|W_0|W}}$, et comme $W_{0S}^{u}\subset W_S^{u}$, on a bien $i_\infty\in W_S^{u}$.

% @ middle, Corollaire and alteration
\uline{Corollaire} 4.4.12. Supposons que $W$ \uline{et} $W_0$ satisfassent : $\mathrm{Loc}(8\ \mathrm{bis})$. Alors on a $D_{W_0}\Longrightarrow A$\unsure{.} (donc \deleted{$D_{W_0}\Longleftrightarrow A$)}\inserted{les huit conditions $A,B,C,D,A_{W_0},B_{W_0},C_{W_0},D_{W_0}$ sont \unsure{équivalentes}}).

% @ middle, explanation
Cela signifie en effet que

% @ lower middle, displayed implication
\[
 A(\mathrm{Parf}_{W_0})\Longrightarrow A(\mathrm{Fibf}_W),
\]

% @ below implication
et c'est \inserted{encore} justiciable du lemme 4.4.9.

% @ bottom paragraph
En effet, la factorisation ci-dessus d'une flèche $f$ ne dépend pas du choix d'un localiseur, et $p_\infty$ est $W$-parfait pour \uline{tout} localiseur satisfaisant $\mathrm{Loc}(8\ \mathrm{bis})$, donc en particulier pour $W_0$. Et la flèche $j_\infty$ est dans $W_S^u$, pour \unsure{tt} localiseur satisfaisant $\mathrm{Loc}(8\mathrm{bis})$, donc pour $W_0$.
% archive sheet number: 332
\scanpage{334}
% @ top centre
\mspage{58}

% @ top paragraph, lines 1–4
Je vais récapituler les résultats obtenus dans cette section en un théorème, qui \unsure{constituera|constitue} le résultat principal du présent \S{} sur \deleted{\unsure{l'}}les axiomes de $W$-\unsure{fibrations|fibration}.

% @ upper middle, theorem, lines 1–6
\uline{Théorème} 4.4.13. Soit $W$ un localiseur fondamental dans $\mathrm{Cat}$, \deleted{$W_0\subset W$}\inserted{\deleted{\unsure{on a les}} Pour \deleted{\unsure{tout}}} un localiseur \inserted{$W_0$} contenu dans $W$, on \deleted{\unsure{trouve}}\inserted{a explicité} un ensemble de \inserted{quatre} conditions sur \unsure{la|le} paire $W_0,W$, qui pour $W$ fixé, dépendent de $W_0$ seul, notées

% @ centre, four conditions
\[
 A_{W_0}, B_{W_0}, C_{W_0}, D_{W_0},\tag*{(4.4.13.1)}
\]
% @ centre, connecting line
et définies \deleted{\unsure{comme}} respectivement comme

% @ centre to lower right, two lines
\[
\begin{aligned}
 &A(\unsure{\mathcal F|\mathscr F}_{W_0}),B(\mathcal F_{\unsure{W|W_0}}),C(\mathcal F_{W_0}),\quad\text{et}\\
 &A(\mathrm{Parf}_{W_0})\ (\Longleftrightarrow B(\mathrm{Parf}_{W_0})\Longleftrightarrow C(\mathrm{Parf}_{W_0})\unsure{)}, 
\end{aligned}
\]
% @ lower middle, reference and added line
avec les notations $A(\cdot),B(\cdot),C(\cdot)$ introduites dans (4.4.1). \deleted{\unsure{$W$}}\inserted{Faisant $W_0=W$, on \deleted{\unsure{a}}\unsure{pose}} \inserted{on trouve les conditions sur $W$}
% @ lower middle, added conditions
\inserted{\[
 A,\ B,\ C,\ D \tag*{(4.4.13.2)}
\]}
% @ lower right, continuation of insertion
\inserted{avec $A=A_W$, $B=B_W$, $C=C_W$, $D=D_W$.}
% @ lower left, beside brace at 4.4.13.2
% slanted
\margin{dont les trois premières ne sont autres que $\mathrm{Loc}(7\unsure{a|*})$, $\mathrm{Loc}(7\mathrm b)$, $\mathrm{Loc}.(7\mathrm c)$ \unsure{(cf} 4.3) pg 38. $\big\}$}

% @ lower quarter, circled numeral
\textcircled{1} On a pour ces conditions les \deleted{deux \unsure{sens} \unsure{ci-|ces}} implications suivantes.

% @ lower quarter, subparagraph and struck preliminary version
a) \deleted{\unsure{On}}\unsure{.|,} les implications \deleted{a) Si $W'_0\subset W_0$ est un localiseur fond. contenu dans $W_0$, on a}
\deleted{\[
 A_{\unsure{W_0|W'_0}}\Longrightarrow A_{W'_0},\ B_{W_0}\unsure{\Rightarrow|=}
\]}
% @ foot, implication chain
\[
 A_{W_0}\Longrightarrow B_{W_0}\unsure{\mathrel{\rule{1.3em}{.5pt}}|\Longrightarrow} C_{W_0}\Longrightarrow D_{W_0}
 \tag*{\deleted{\unsure{(4.4.13.2)}}(4.4.13.3)}
\]
% @ foot, unfinished line
(d'où en particulier
% archive sheet number: 333
\scanpage{335}
% @ top centre
\mspage{59}

% @ top, continuation of parenthesis
\[
 A_W\Longrightarrow B_W\Longrightarrow C_W\Longrightarrow D_W\deleted{\unsure{.}}\big).\tag*{(4.4.13.4)}
\]
% @ top, b), two lines
b) Si $W'_0\subset W_0$ est un localiseur \inserted{fond.} contenu dans $W_0$, on a (tautologiquement)

% @ upper quarter, three implications
\[
 A_{W_0}\Rightarrow A_{W'_0},\ B_{W_0}\Rightarrow B_{W'_0},\ C_{W_0}\Rightarrow C_{W'_0}.
\]
% @ above and beside diagram
On \deleted{\unsure{a}} trouve donc par a) et b) \inserted{(y faisant \deleted{\unsure{possiblement|provisoirement}} $(W,W_0)$ comme $(W_0,W'_0)$)}, le rectangle d'implications

% @ centre, two columns, two side braces and two underbraces at the foot
\diagram{\begin{tikzcd}[ampersand replacement=\&,row sep=1.1em,column sep=1em]
 \& A\arrow[r,Rightarrow,"{\deleted{\unsure{\leftarrow|\cdot}}}" description,very near start]\arrow[d,Rightarrow] \& A_{W_0}\arrow[d,Rightarrow] \& {}\\
 \unsure{\mathcal F|\mathscr F}_W,\mathrm{Fibf}_W\,\left\{\vphantom{\rule{0pt}{2.8em}}\right. \& B\arrow[r,Rightarrow]\arrow[d,Rightarrow] \& B_{W_0}\arrow[d,Rightarrow]\& \left.\vphantom{\rule{0pt}{2.8em}}\right\}\,\mathcal F_{W_0},\mathrm{Fibf}_{W_0}\\
 (4.4.13.4) \& C\arrow[r,Rightarrow]\arrow[d,Rightarrow] \& C_{W_0}\arrow[d,Rightarrow]\& {}\\
 \& \underbrace{D}_{\mathrm{Parf}_W}\arrow[r,Rightarrow]\& \underbrace{D_{W_0}}_{\mathrm{Parf}_{W_0}}\& {,}
\end{tikzcd}}

% @ below diagram, paragraph, lines 1–5
où on a \deleted{\unsure{noté déjà}} rappelé, à côté \unsure{des} accolades, quels sont \deleted{\unsure{les}} l'ensemble\inserted{, ou les ensembles,} $\mathcal F\subset\mathrm{Fl}(\mathrm{Cat})$ qui définissent les conditions envisagées \deleted{\unsure{par}} (comme l'une des conditions $A(\mathcal F), B(\mathcal F), C(\mathcal F)$ ou $D(\mathcal F)=A(\mathcal F)$ dans \unsure{le cas|les cas} où les trois conditions sont équivalentes).

% @ lower quarter, comparative sentences
La plus forte des huit conditions envisagées est $D_{W_0}$, la plus forte $A$. \deleted{\unsure{Le}} (Et pour $W_0$ \unsure{variable}, on trouve la condition la plus faible de toutes en prenant $D_{W_\infty}$, où $W_\infty$ est le plus petit des localiseurs fondamentaux dans $\mathrm{Cat}$.)

\figure{\begin{tikzpicture}\draw (0,0) -- (.70,.12) -- (.70,.60); \draw (.65,.16) -- (.65,.63); \draw (.71,.66) -- (1.23,.66);\end{tikzpicture}}

% @ bottom, c)
c) On a de plus, pour la condition la plus faible $D_{W_0}$, l'implication
% @ bottom edge, display
\[
 D_{W_0}\Longrightarrow F,
\]
% archive sheet number: 334
\scanpage{336}
% @ top centre
\mspage{60}

% @ top three lines
où $F=\mathrm{Loc}(7\mathrm F)$ est la condition que $\mathrm{Cat}$ possède suffisamment de $W$-fibrations, i.e. que c'est une catégorie à fibrations pour $(W,\mathrm{Fib}_W)$.

% @ upper middle, circled 2 and following six lines
\textcircled{2} Si $W$ et $W_0$ satisfont à $\mathrm{Loc}(8\mathrm{bis})$ (axiome des \unsure{limites} dénombrables), alors \deleted{\unsure{toutes}} $D_{W_0}\Longrightarrow A$, i.e les 8 conditions envisagées sont équivalentes. Donc si $W$ satisfait $\mathrm{Loc}(8\mathrm{bis})$, \deleted{\unsure{prenant pour $W_0$}}\inserted{considérant} le plus petit loc fondamental \inserted{$W_L$} satisfaisant à $\mathrm{Loc}(8\ \mathrm{bis})$ (de sorte que $W\supset W_L$), on voit que les huit conditions

% @ middle, eight conditions
\[
 A,B,C,D,\quad A_{W_L},B_{W_L},C_{W_L},D_{W_L}
\]
% @ middle left beside list
% slanted
\margin{\deleted{\unsure{NB}} la forme la plus faible, $D_{W_L}$, de ces huit \uline{conditions, est} la relation $\boxed{\mathrm{Parf}_{W_L}\cap W\subset W^{\unsure{u|n}}}$}
% @ below list
sont équivalentes. En particulier, sous la hypothèse (\deleted{\unsure{fort}} anodine) sur $W$, les conditions \deleted{\unsure{$\mathrm{Loc}(7$}}\inserted{$\mathrm{Loc}(7)=$} $\mathrm{Loc}(7\mathrm a)$, $\mathrm{Loc}(7\mathrm b)$ et $\mathrm{Loc}(7\mathrm c)$ sont équivalentes, et elles équivalent aussi à la condition $D\unsure{=|\Leftarrow}A(\mathrm{Parf}_W)$ i.e. à

% @ lower middle, inclusion
\[
 \mathrm{Parf}_W\subset\mathrm{Fib}_W,\quad\text{(\deleted{\unsure{cas}} particulier de \inserted{$\mathrm{Loc}(7\mathrm a)$}),}
\]
% @ lower quarter
et à ses deux variantes \deleted{\unsure{équivalentes}} d'apparence plus faible $B(\mathrm{Parf}_W)$, $C(\mathrm{Parf}_W)$ (cas particuliers de $\mathrm{Loc}(7\mathrm b)$ et de $\mathrm{Loc}(7\mathrm c)$ respectivement).

% @ bottom two lines
\textcircled{3} Supposons que $W_0$ \unsure{satisfasse} à $\mathrm{Loc}(8\mathrm{bis})$ (sans hypothèse particulière sur $W$). Alors
% archive sheet number: 335
\scanpage{337}
% @ top centre
\mspage{61}

% @ top three lines
les quatre conditions $A_{W_0},B_{W_0},C_{W_0},D_{W_0}$ sont équivalentes, et le diagramme d'implications de \deleted{\unsure{2}}1\textsuperscript{o}) (p.59) se réduit donc au diagramme linéaire

% @ upper centre, chain
\[
 A\Rightarrow B\Longrightarrow C\Longrightarrow D\Longrightarrow D_{W_0}\ (\Leftrightarrow A_{W_0}\ldots),
\]
% @ below chain
dont le premier terme \deleted{est la relation} est $\mathrm{Loc}(7\unsure{\mathrm a|*})=\mathrm{Loc}(7)$, et le dernier équivaut à l'inclusion

% @ first inclusion, upper third
\[
 \deleted{\unsure{\mathrm{Fi}}}\unsure{\mathcal F|\mathscr F}_{W_0}\subset\mathrm{Fib}_W,
\]
% @ immediately below first inclusion
\inserted{cas particulier de $A=\mathrm{Loc}(7\mathrm a)$,} ou encore à l'inclusion (en apparence plus forte)

% @ second inclusion
\[
 \mathrm{Fibf}_{W_0}\subset\mathrm{Fib}_W
\]
% @ above third inclusion
(cas particulier de la forme équivalente
% @ third inclusion
\[
 \mathrm{Fibf}_W\subset\mathrm{Fib}_W
\]
% @ centre, end of parenthesis
de l'axiome $\mathrm{Loc}(6\unsure{\mathrm c|\mathrm a})$).

% @ centre, circled 4
\textcircled{4} Enfin, \deleted{\unsure{signalons à part}} les implications plus techniques conditionnelles suivantes

% @ lower middle, a)
a) Si $W_0$ satisfait à $\mathrm{Loc}(7\mathrm F)$, alors les conditions $A_{W_0},B_{W_0},C_{W_0}$ sont équivalentes, donc le diagramme (\unsure{4.4413.4|4.4.13.4}) se réduit au diagramme

% @ lower quarter, two-column diagram
\diagram{\begin{tikzcd}[ampersand replacement=\&,row sep=small]
 A\arrow[d,Rightarrow]\& {}\& {}\\
 B\arrow[d,Rightarrow]\& {}\& {}\\
 C\arrow[d,Rightarrow]\arrow[rr,Rightarrow]\& {}\& C_{W_0}\arrow[d,Rightarrow]\\
 \unsure{D|b}\& \unsure{\mathrel{=\joinrel=}|\Rightarrow}\& D_{W_0}\& {.}
\end{tikzcd}}
% @ bottom edge, b) and correction at right
\deleted{\unsure{a)}}b) Si $W$ est fortement saturé (\deleted{condition \unsure{sur}}\inserted{axiome $\mathrm{Loc}(1\unsure{\mathrm{ter}|\mathrm{bis}})$}), \unsure{quel|qui}
% archive sheet number: 336
\scanpage{338}
% @ top centre
\mspage{62}

% @ top paragraph, with interlinear revision below Der 2
est satisfaite dans la grande majorité des cas, et est même impératif pour construire un prédérivateur satisfaisant à l'axiome de localisation $\mathrm{Der}\ 2$), \deleted{alors on a $B\Rightarrow A$}\inserted{et satisfait à $\mathrm{Loc}(3\mathrm{bis})$ (axiome des \deleted{\unsure{cofibrations}} localisations \unsure{fortes|finies}, cf \unsure{XI|\S{} 1})}, alors on a $B\Rightarrow A$, \deleted{\unsure{donc}} donc

% @ upper quarter, equivalence
\[
 A\Longleftrightarrow B,
\]
% @ line beneath equivalence
i.e. $\mathrm{Loc}(7\unsure{\mathrm c|\mathrm a})=\mathrm{Loc}(7)$ \deleted{\unsure{la|le}} équivaut à $\mathrm{Loc}(7\mathrm b)$.

% @ upper third, struck introductory line
\deleted{Je résume les implications}

% @ middle, proof through chapter citation
\uline{Dém} \deleted{\unsure{Toutes}} Tout ce qui est énoncé a été vu dans les pages précédentes, à la seule exception de $4^\circ$, c. Mais cette implication $B\Rightarrow A$ pour \deleted{\unsure{$A$}}$W$ fort. saturé est prouvé dans $\underline{\overline{\unsure{\mathrm{VII}|\mathrm{XII}}}}$ (prop. 1, p.54).

% @ middle and lower middle, recollection of proof
Je rappelle que la démonstration utilise de façon essentielle le formalisme du \deleted{catégories} prédérivateur $\mathrm{HOT}_W$ (construit et étudié surtout dans $\underline{\overline{\mathrm{VI}}}$), et \deleted{\unsure{entière}}\inserted{la théorie} \unsure{du foncteur|des foncteurs} $\unsure{d_!|d!|f_!}$, et notamment le résultat (valable si $W$ satisfait à $\mathrm{Loc}(\inserted{\unsure{3\mathrm{bis}|8\mathrm{bis}},}7\mathrm b)$) que pour une flèche $f:S'\unsure{\mathrel{-}|\to} S$ dans $W$, la flèche correspondante

% @ lower quarter, displayed map
\[
 f^{\unsure{*|\prime}}_{\mathrm{lc}}:\mathrm{Hot}^{\mathrm{lc}}_W(S)\longrightarrow\mathrm{Hot}^{\mathrm{lc}}_W(\unsure{S'|B'})
\]
% @ bottom, final three lines
est une équivalence de catégories. Donc l'implication $B\Rightarrow A$, elle aussi, est un résultat qui n'a rien d'évident.
% archive sheet number: 337
\scanpage{339}
% @ top centre
\mspage{63}

% @ above section number
\deleted{\uline{\unsure{Commentaire|Corollaire}} 4.4.14.}

% @ upper paragraph, four lines
4.4.14. Je résume les implications obtenues jusqu'à présent \deleted{dans le \unsure{diag}}, y compris les implications conditionnelles, dans le diagramme récapitulatif suivant :

% @ centre, recapitulation diagram
\diagram{\begin{tikzcd}[ampersand replacement=\&,column sep=3.6em,row sep=2.1em]
 |[circle,draw,inner sep=2pt]| A\arrow[r,Rightarrow]\arrow[d,Rightarrow] \& A_{W_0}\arrow[d,Rightarrow] \& {}\\
 B\arrow[r,Rightarrow]\arrow[d,Rightarrow]\arrow[u,bend left=57,dashed,"{W(1\mathrm{ter},3\mathrm{bis})}" description,"{\deleted{\unsure{W(1\mathrm{ter})}}}" near end] \& B_{W_0}\arrow[d,Rightarrow] \& {}\\
 C\arrow[r,Rightarrow]\arrow[d,Rightarrow] \& C_{W_0}\arrow[d,Rightarrow]\arrow[uu,bend right=42,dashed,"{W_0(\unsure{7|7\mathrm F})}" description] \& {}\\
 D\arrow[r,Rightarrow]\arrow[uuu,bend left=65,dashed,preaction={draw,solid,decorate,decoration={zigzag,segment length=3pt,amplitude=2.5pt}},"{\deleted{\unsure{W(8)}}}" description]\arrow[dr,dashed,"{W(8\mathrm{bis})}"'] \& D_{W_0}\arrow[r,Rightarrow]\arrow[uuu,bend right=54,dashed,"{W_0(8\mathrm{bis})}" description]\arrow[d,Rightarrow,"{\substack{\deleted{\unsure{W_0(8\mathrm{bis})}}\\ W(8\mathrm{bis})+W_0(8\mathrm{bis})}}"]\arrow[d,shift right=1.1,densely dotted] \& F\\
 \& |[circle,draw,inner sep=2pt]| A\& {,}
\end{tikzcd}}
% @ slanting above arrow to F
% slanted
\margin{NB on verra sur l'exemple de $\unsure{\mathbb W_0|W_0}$ que cette implication est \uline{stricte}, même sous l'hypothèse $\mathrm{Loc}(1\ \mathrm{ter},2\ \mathrm{bis},4,5,6\mathrm a,8,9)$}
% @ below diagram, four lines
où j'ai cerclé la condition la plus forte et la plus importante de toutes, $A=\mathrm{Loc}(7)$, \unsure{qui|que} \unsure{apparait|apparaît} aux deux extrémités opposées du diagramme.

% @ bottom paragraph, five lines
Je voudrais encore faire un commentaire sur les ingrédients des résultats prouvés ici, et notamment de l'équivalence des huit conditions $A$ à $D_{W_0}$, quand $W$ et $W_0$ satisfont $\mathrm{Loc}(8\mathrm{bis})$ :
% archive sheet number: 338
\scanpage{340}
% @ top centre
\mspage{64}

% @ top, item 1 and interline
1) La théorie des \uline{catégories de chemins} $\underline{\mathrm{Ch}}(X)$ ($\underline{\overline{\mathrm{VII}}}$), \inserted{s'appuyant sur celle des \uline{morphismes lisses}, \uline{propres}\unsure{,} \uline{parfaits}.}

% @ top, item 2
2) Le lemme de K.S. Brown (3.2.1 \unsure{pg|pag|page} 20)

% @ upper third, item 3
3) Pour l'implication $B\Rightarrow A$ i.e. $\mathrm{Loc}(7\mathrm b)\Rightarrow \mathrm{Loc}\,7$, (sous l'hypothèse $W(1\ \mathrm{ter},\deleted{\unsure{8\mathrm{bis}}}3\ \mathrm{bis})$), j'utilise la théorie du prédérivateur $\mathrm{HOT}_W$ (VI), et plus particulièrement celle des foncteurs $f_{\unsure{!|.}}$ (définis sous la seule hypothèse que \deleted{$W$ satisfasse} \deleted{\unsure{$\mathrm{Loc}$}}$W(3\ \mathrm{bis})$, renforçant la condition automatique $W(3)$ (faisant partie de la définition d'un loc fond.)

% @ centre and lower middle, Brown and open question
Il est clair que je ne serais pas arrivé à y voir clair, sans le travail fondamental de K.S. Brown. Il reste pourtant encore des questions ouvertes. Ainsi, si on ne suppose pas que $W$ satisfasse à $\mathrm{Loc}(8\ \mathrm{bis})$, mais on \unsure{borne|se borne|tient} à $\mathrm{Loc}(3\ \mathrm{bis})$ (et je n'ai pas eu l'occasion encore de travailler avec un localiseur qui n'y satisfasse \ldots), est-il vrai néanmoins que l'on a \inserted{encore} l'implication

% @ lower quarter, question's implication
\[
 D\Longrightarrow A,
\]
% @ foot, question continues
de façon que les conditions $A,B,C,D$ sont équivalentes (et en particulier $\mathrm{Loc}(7\mathrm c)$ implique déjà $\mathrm{Loc}(7)$ dans \deleted{\unsure{toute}} sa force) ? Il est clair dans tout cela que
% archive sheet number: 339
\scanpage{341}
% @ top centre, overwritten numeral
\mspage{6\deleted{\unsure{4}}5}

% @ top paragraph, first two lines
je ne me suis pas suffisamment attaché à développer des exemples et des \unsure{contr'exemples|contre'exemples}.

% @ upper third, observations and hypothetical implication
Sans compter qu'il me semble que tous les loc. fond. \inserted{sur $\mathrm{Cat}$} qu'on rencontre ``en pratique'' satisfont déjà à $\mathrm{Loc}(8)$ et a fortiori à $\mathrm{Loc}(8\mathrm{bis})$, donc pour eux on a \deleted{\unsure{bien}} l'implication escomptée. La \unsure{$=$|même} remarque s'applique à l'implication hypothétique

% @ centre, question mark over arrow
\[
 D_{W_0}\overset{?}{\Longrightarrow}A,
\]
% @ centre, qualification
sans hypothèse $\mathrm{Loc}(8\ \mathrm{bis})$ sur $W_0$ et sur $W$.

% @ lower middle through foot, candidate and next question
\deleted{Il faudrait \unsure{exp}} Un candidat évident pour un contre-exemple déjà à la première question serait $W_\infty={}$ plus petit localiseur fondamental sur $\mathrm{Cat}$, duquel il est presque évident qu'il ne satisfait pas à $\mathrm{Loc}(8\deleted{\unsure{\mathrm{bis}}})$ \unsure{ni} même $\mathrm{Loc}(8\mathrm{bis})$. Mais il n'y a aucune raison qu'il satisfasse \inserted{non plus} à $F$, \deleted{donc} de telle façon que (s'il n'y satisfait pas\unsure{\deleted{)}|)}, donc non plus à $D$) l'implication \inserted{$D\Rightarrow A$} est bel et bien valable. Cela nous amène à la question, plus substantielle, si on n'aurait pas \deleted{\unsure{cela}} l'implication
% archive sheet number: 340
\scanpage{342}
% @ top centre
\mspage{66}

% @ top, displayed question
\[
 F\overset{?}{\Longrightarrow} A,
\]
% @ upper left beside question
% slanted
\margin{C'est faux, cf plus bas l'implication $\unsure{\mathrm{Doc}|\mathrm{Loc}}(7D_{W_\infty})\Longrightarrow \mathrm{Loc}(7\mathrm F)$ est \uline{stricte}.}
% @ top paragraph below question
du moins \unsure{ss|sous} des conditions convenables assez anodines sur $W$, genre $\mathrm{Loc}(1\ \mathrm{ter},3\ \mathrm{bis},8\ \mathrm{bis})$.

% @ upper to middle paragraph
C'est sans doute là la question la plus substantielle \inserted{et la plus intéressante} qui est soulevée par le diagramme de la page 63, et qui ne fait plus appel à l'introduction d'un localiseur fondamental $W_0$ auxiliaire. \deleted{Après}\inserted{Vu} le lemme de Brown, il ne me semble pas absolument exclu qu'un tel miracle ait bel et bien lieu !

% @ centre, interlined heading and section opening
4.4.15. \inserted{(Digression sur l'axiome $\mathrm{Loc}(6)$ des cofibrations)} Cela attire mon attention sur la question en quelque sorte duale : Si $\mathrm{Cat}$ a assez de $W$-\uline{co}fibrations, i.e si \unsure{muni|munie} de $(W,\mathrm{Cof}_W)$ c'est une catégorie à cofibrations, cela implique-t-il déjà l'axiome $\mathrm{Loc}(6\mathrm a)$\unsure{$*$|.|$?$} (qui avait été dégagé, essentiellement, pour assurer qu'il en soit bien ainsi). La proposition 3.6. \inserted{(p.25)} nous dit qu'il revient au \unsure{$=$|même} de prouver que, si $W$ satisfait $\mathrm{Cof}$, alors les carrés de M.V. ouverts sont
% archive sheet number: 341
\scanpage{343}
% @ top centre
\mspage{67}

% @ top paragraph, lines 1–4
\deleted{déjà} $W$-\unsure{cocartésiens|cocartésien}. \deleted{\unsure{Pourtant}} Mais il me semble à présent que nous pouvons le prouver, \inserted{\unsure{i.e.}} sous la seule condition supplémentaire \deleted{\unsure{$W$}} $\mathrm{Loc}(3\mathrm{bis})^{*}$,
% @ upper left opposite asterisk, three diagonal lines
% slanted
\margin{et peut-être, de plus, $\mathrm{Loc}(1\ \mathrm{ter})$ $W$ fortement saturé}
% @ upper third, continuation to diagram pair
laquelle assure l'existence \unsure{des foncteurs|du foncteur} $f_!$, et en particulier la construction du \deleted{\unsure{carré}} $\mathrm{Holim}$ d'un diagramme

% @ upper third, two forks with interposed comme
\diagram{\begin{tikzcd}[ampersand replacement=\&,row sep=0pt,column sep=small]
 \& \& X_1 \& \& \& X_1\\
 \& X_0\arrow[ur]\arrow[dr]\& {} \& \text{comme}\quad\displaystyle\int \& X_0\arrow[ur]\arrow[dr]\& {}\\
 \deleted{\unsure{\int|f|\text{Or}}},\& \& X_{\unsure{2|\sim}}\& \& \& X_2.
\end{tikzcd}}
% @ centre, statement before open square
Donc dire qu'un diagramme de MV ouvert

% @ centre, open square in diamond position
\diagram{\begin{tikzcd}[ampersand replacement=\&,row sep=0pt,column sep=small]
 \& X_1\arrow[dr]\& {}\\
 X_0\arrow[ur]\arrow[dr]\& \& X\\
 \& X_2\arrow[ur,dash]\& {}
\end{tikzcd}}
% @ centre to lower, canonical map
est $W$-cocartésien, doit revenir à dire que la flèche canonique

% @ lower middle, integral of fork and map
\diagram{\begin{tikzcd}[ampersand replacement=\&,row sep=0pt,column sep=small]
 \& \& X_1\& {}\\
 \displaystyle\int\& X_0\arrow[ur]\arrow[dr]\& {}\arrow[r]\& X\\
 \& \& X_2\& {}
\end{tikzcd}}
% @ lower third, qualification and references
est dans $W$; \deleted{\unsure{du}} \unsure{moins|disons} si $W$ est supposé fortement saturé. Or cela a été bel et bien vu dans $\underline{\overline{\mathrm{XII}}}$ 3.4. Th.4 et son \uline{corollaire} 3, p.149). Dans loc cit. j'énonce les conclusions sous l'hypothèse que \inserted{\unsure{j'appelais}} $W(\unsure{A|7}\mathrm{bis})$ (l'actuel $\mathrm{Loc}(6)$), mais je viens de vérifier que pour la partie \uline{directe}, donc aussi pour ses cor. 1 à 3, on n'a pas besoin
% archive sheet number: 342
\scanpage{344}
% @ top centre
\mspage{\unsure{58|68}}

% @ top two lines
de cette hypothèse (comme je l'indique en annotations marginales).

% @ upper third, numbered paragraph
4.16. Le résultat que je viens de prouver (sans vérifier \unsure{bien sûr|de visu sur} la démonstration) me paraît aussi délicat, qu'il est satisfaisant. Cela \deleted{\unsure{donne à}} rend la ``conjecture'' \unsure{hasardeuse|hasardée|hésitante} de 4.4.14, savoir que

% @ upper middle, double implication
\[
 F\overset{?}{\Rightarrow}A\quad\text{donc}\quad F\Leftrightarrow A
\]
% @ beneath conjecture, four lines
en supposant que $W$ satisfait $\mathrm{Loc}(1\ \mathrm{ter},3\ \mathrm{bis},8\ \mathrm{bis})$, nettement plus engageante. \unsure{Car|Or} on commence à s'habituer aux miracles les plus époustouflants \ldots

% @ middle, short separating stroke
\figure{\begin{tikzpicture}\draw (0,0) -- (.28,-.3);\end{tikzpicture}}

% @ lower middle, counterexample and insertion
Pourtant, je viens de vérifier que cette conjecture est fausse, sur l'exemple de $\unsure{\mathbb W_0|W_0}\inserted{=\{f\in\mathrm{Fl}(\mathrm{Cat})\mid\pi_0(f)\ \mathrm{iso}\}}$, qui satisfait à $F$, mais pas à $A$, ni même à $D_{W_\infty}$. Ainsi l'implication

% @ lower middle, single implication
\[
 D_{W_\infty}\Longrightarrow F
\]
% @ before brace
et a fortiori les implications
% @ lower quarter, brace enclosing two lines
\[
 \left\{\begin{aligned}
 &\mathrm{Loc}(7\mathrm a)\Rightarrow F,\ \mathrm{Loc}(7\mathrm b)\Rightarrow F,\ \mathrm{Loc}(7\mathrm c)\Rightarrow F\\
 &D_{W_0}\Longrightarrow F
 \end{aligned}\right.
\]
% @ foot, restriction
sont \uline{strictes}, même en se \unsure{limitant|restreignant} à des localiseurs fondamentaux $W$ satisfaisant (comme $\mathbb W_0$) à
% @ bottom display
\[
 \mathrm{Loc}(1\ \mathrm{ter},2,\unsure{B|3|2}\mathrm{bis},4,5,6\mathrm a,8,9).
\]
% @ lower right beside the list, vertical dividing stroke
\margin{\begin{tabular}{|p{1.85cm}@{}}\uline{NB} $\mathbb W_0$ ne satisfait \uline{pas} $\mathrm{Loc}(6\mathrm b)$, ni $\mathrm{Loc}(7D_{W_\infty})$ cf plus bas) par contre $\mathrm{Loc}(\mathrm F)$.\end{tabular}}
% archive sheet number: 343
\scanpage{345}
% @ top centre
\mspage{69}

% @ top through parenthetical NB
Que $\unsure{\mathbb W_0|W_0}$ satisfasse à cet ensemble d'axiomes se vérifie \deleted{de façon} par \unsure{AQT|A.Q.T.}. (\uline{NB} Qu'il satisfasse à $\mathrm{Loc}(1\ \mathrm{ter},2,3\mathrm{bis})$ est déjà conséquence du fait que $\mathbb W_0$ est de la forme $\deleted{\unsure{\mathbb D}}W_{\mathbb D}$, où $\mathbb D$ ($=\mathbb D_{\mathbb W_0}$ en l'occurrence) est un prédérivateur satisfaisant $\mathrm{Der}(3,4,5)$, ce qui est bien le cas pour $\mathbb D_{\mathbb W_0}$, qui est un vrai \uline{dérivateur}.)

% @ centre and lower half, factorization
Il faut vérifier que tout $f$ dans $\mathrm{Cat}$ se factorise en $f=pi$, avec $i\in\mathbb W_0$, $p$ une $\mathbb W_0$-fibration. Or si $f:X\longrightarrow Y$, $I=\pi_0(X)$, on prend $\overline X=\coprod_{\unsure{d|\alpha}\in I}Y_{f(d)}$, \unsure{d'où} une factorisation $f=pi$, où $\deleted{\unsure{i\in}}\pi_0(i)$ iso i.e. $i\in\mathbb W_0$, et où $p:\overline X\longrightarrow\deleted{\unsure{X|Y}}$ est, au dessus de la composante connexe $Y_j$ de $Y$, donné par le produit $\overline X\mid Y_j\simeq Y_j\times\underbrace{\pi_0(f)^{-1}(j)}_{\text{cat. discrète}}$, \inserted{$p$} \deleted{\unsure{$q$}} est bien une $\mathbb W_0$-fibration, car grâce à $\mathrm{Loc}(4)$ il suffit de le vérifier sur chaque $Y_j$, et là c'est déduit d'une flèche $I_j\unsure{\longleftrightarrow|\longrightarrow}e$ (une $\mathbb W_0$-fibration) par chgt de base $Y_j\longrightarrow e$.

% @ foot, beginning of next argument
\deleted{\unsure{on|un}}\unsure{/|,} D'autre part, que $\mathbb W_0$ ne \unsure{satisfait|satisfasse} \uline{pas} à $D_{\unsure{\mathbb W_\infty|W_\infty}}$ \deleted{\unsure{est}} a déjà été vu ($\overline{\mathrm I}$, p.87),
% archive sheet number: 344
\scanpage{346}
% @ top centre
\mspage{70}

% @ top left sloping note
% slanted
\margin{Autre façon de le voir : on sait \inserted{($\underline{\overline{\mathrm{XII}}}$)} que $W(5,7D_{W_\infty})\Longrightarrow W\subset\unsure{\mathbb W_\infty|W_\infty}$, or on a $\mathbb W_\infty\subset\mathbb W_0$, et l'inclusion est \uline{stricte} !}
% @ top paragraph, four lines
réalisant l'application de \unsure{revêtement} $f:S^1\longrightarrow S^1$ dans $\mathrm{Cat}$ par un morphisme $W_\infty$-parfait, (en fait, par une \inserted{$\mathrm{Cat}$} fibration localement triviale), chose très facile.

% @ upper third, separating stroke
\figure{\begin{tikzpicture}\draw (0,0) -- (.37,-.37);\end{tikzpicture}}

% @ middle, five lines before diagram
Je pense qu'à présent \inserted{que l'ensemble} des différentes variantes de l'axiome des $W$-fibrations est à peu près compris, à cela près qu'il n'est toujours pas clair, \inserted{\uline{en dehors de l'axiome} $\mathrm{Loc}(8\mathrm{bis})$,} si les implications du diagramme

% @ central rectangle
\diagram{\begin{tikzcd}[ampersand replacement=\&,row sep=small]
 A\arrow[r,Rightarrow]\arrow[d,Rightarrow]\& A_{W_0}\arrow[d,Rightarrow]\\
 B\arrow[r,Rightarrow]\arrow[d,Rightarrow]\& B_{W_0}\arrow[d,Rightarrow]\\
 C\arrow[r,dash,double]\arrow[d,Rightarrow]\& C_{W_0}\arrow[d,Rightarrow]\\
 D\arrow[r,Rightarrow]\& D_{W_0}
\end{tikzcd}}
% @ lower third, completion
sont strictes ou non. Mais il me semble que tous les localis\deleted{\unsure{at}}eurs fondamentaux intéressants rencontrés jusqu'à présent satisfont $\mathrm{Loc}(8)$, et même $\mathrm{Loc}(1\ \mathrm{ter},2\deleted{\unsure{\mathrm{bis}}},3\ \mathrm{bis},4,5,6\mathrm a,8,9)$ (peut-être même $\mathrm{Loc}(7\mathrm F)$ -- il \margin{?} faudrait que je teste sur les $\mathbb W_{\unsure{n|m}}$), donc la perplexité non résolue me paraît sans \unsure{grande conséquence|grandes conséquences}.
% archive sheet number: 345
\scanpage{347}
% @ top centre, manuscript pagination
\mspage{71}

% @ upper third, crossed block, lines 1–7
\deleted{Pour terminer, je voudrais donner encore une reformulation de $A=\mathrm{Loc}(7\mathrm{a})$, très voisine de celle donnée dans le cor. 4.4.5 (page 45), où on prouve que $A$ équivaut à $C$, \uline{plus} un ``petit \unsure{chouïa|chouia}'' (pas si petit que \unsure{sa|ça} !). Ici, on va remplacer $C$ par $D$, voire par $D_{W_0}$.}

% @ heading, lines 8–9
5. \uline{Axiome des limites} $\mathrm{Loc}(8)$, \uline{et l'axiome $\mathrm{Loc}$}$(9)$ \uline{d'accessibilité}.

% @ line 10
C'est l'axiome :

% @ line 11
\uline{Loc} 8 : $W$ est stable par $\underset{\longrightarrow}{\lim}$ filtrantes.

% @ line 12
Une variante affaiblie utile est la suivante

% @ line 13
\uline{Loc} 8bis : $W$ est stable par $\underset{\longrightarrow}{\lim}$ dénombrables.

% @ lines 14–15
Explicitement, si on a \unsure{deux} systèmes inductifs filtrants dans $\deleted{\unsure{X|x}}\mathrm{Cat}$

% @ lines 16–17
\[
\begin{aligned}
\underline X\unsure{\cdot|.}&=(X_i)_{i\in I}\\
\underline Y&=(Y_i)_{i\in I},
\end{aligned}
\]

% @ line 18
et une flèche

% @ line 19
\[
\underline f=(f_i)_{i\in I}:\underline X\longrightarrow\underline Y,
\]

% @ lines 20–23
donnée par des $f_i:X_i\longrightarrow Y_i$, \deleted{alors Loc 8 nous dit que si} \inserted{posant} $X=\underset{\longrightarrow}{\lim}X_i$, $Y=\underset{\longrightarrow}{\lim}Y_i$, $f=\underset{\longrightarrow}{\lim}f_i:X\mathbin{\text{---}}Y$, Loc 8 nous dit que si les $f_i\in W$, alors $f\in W$. Ici $I$ est

% archive sheet number: 346
